\documentclass[11pt,a4paper]{article}

% ====== ПАКЕТЫ ======
\usepackage[utf8]{inputenc}
\usepackage[T1,T2A]{fontenc}
\usepackage[english,russian]{babel}
\usepackage{geometry}
\geometry{a4paper, left=1.5cm, right=1.5cm, top=1cm, bottom=1cm}
\usepackage{xcolor}
\usepackage{enumitem}
\usepackage{hyperref}
\usepackage{titlesec}
\usepackage{tabularx}

% ====== ЦВЕТА ======
\definecolor{maincolor}{HTML}{2F9E43}
\definecolor{darktext}{HTML}{1A1A1A}
\definecolor{lighttext}{HTML}{666666}

\hypersetup{colorlinks=true, linkcolor=maincolor, urlcolor=maincolor}
\pagestyle{empty}

% ====== ЗАГОЛОВКИ ======
\titleformat{\section}
  {\Large\bfseries\color{maincolor}}
  {}{0em}{}[\titlerule]
\titlespacing{\section}{0pt}{8pt}{3pt}

% ====== НАСТРОЙКИ ======
\setlist{leftmargin=*, nosep, itemsep=2pt}
\setlength{\parindent}{0pt}

% ====== КОМАНДЫ ======
\newcommand{\cvevent}[4]{%
  \vspace{3pt}
  \noindent
  {\bfseries\large #1} \hfill {\bfseries #4}\\
  {\color{lighttext}\textit{#2}} \hfill {\color{lighttext}\textit{#3}}
  \vspace{1pt}
}

% ====== НАЧАЛО ======
\begin{document}

% ===== ШАПКА =====
\begin{center}
  {\Huge\bfseries\color{darktext} Ольга Зырянова}
  \vspace{2pt}
  
  {\large\color{lighttext} Backend Developer | Computer Science Student}
  \vspace{6pt}
  
  \small
  Москва, Россия \quad
  +7 908 704-04-61 \quad
  \href{mailto:olga.zyryanova14@gmail.com}{olga.zyryanova14@gmail.com} \\
  \href{https://github.com/ssuunneedd}{github.com/ssuunneedd} \quad
  \href{https://codeforces.com/profile/olyazyryanova}{codeforces.com}
\end{center}

\hrule
\vspace{4pt}

% ===== О СЕБЕ =====
\section*{О себе}
Студентка 1-го курса ФКН НИУ ВШЭ с опытом разработки и преподавания. 
Призёр олимпиад по информатике и математике.

% ===== ОБРАЗОВАНИЕ =====
\section{Образование}

\cvevent{НИУ ВШЭ, ФКН}{Прикладная математика и информатика}{2025 -- 2029}{Москва, Россия}
\begin{itemize}
  \item Средний балл: \textbf{8.88/10.0}
\end{itemize}

\cvevent{Школа ЦПМ}{Математика и информатика}{2022 -- 2025}{Москва, Россия}

\cvevent{МБОУ ФМЛ №31}{Основное общее образование}{2018 -- 2022}{Челябинск, Россия}

% ===== ОПЫТ РАБОТЫ =====
\section{Опыт работы}

\cvevent{Преподаватель}{Летняя компьютерная школа}{2025 -- 2026}{Кострома, Россия}
\begin{itemize}
  \item Провела лекции и практики для \textbf{30+ школьников} по алгоритмам и структурам данных
  \item Разработала \textbf{15 учебных заданий}
\end{itemize}

% ===== НАВЫКИ =====
\section{Навыки}

\begin{tabularx}{\textwidth}{@{}lX@{}}
  {\color{lighttext}\textbf{Языки:}} & Python, C++, Kotlin \\
  {\color{lighttext}\textbf{Базы данных:}} & PostgreSQL \\
  {\color{lighttext}\textbf{Data Science:}} & NumPy, Pandas, Spark, Matplotlib, Seaborn \\
  {\color{lighttext}\textbf{Backend:}} & FastAPI, Django \\
  {\color{lighttext}\textbf{Frontend:}} & Vue.js, JavaScript, HTML, CSS \\
  {\color{lighttext}\textbf{Инструменты:}} & Git, Linux, Docker, LaTeX \\
  {\color{lighttext}\textbf{Английский:}} & B2 \\
\end{tabularx}

% ===== ВОЛОНТЁРСТВО =====
\section{Волонтёрство и общественная деятельность}

\cvevent{Благотворцы}{Организатор благотворительной инициативы}{2026 -- наст. время}{Россия}
\begin{itemize}
  \item Организация благотворительных мастер-классов и мероприятий по рукоделию.
  \item Координация работы волонтёров и взаимодействие с участниками проектов.
  \item Привлечение средств и материалов для проведения благотворительных акций.
\end{itemize}

\cvevent{Волонтёрская деятельность}{Участник и организатор мероприятий}{2025 -- наст. время}{Россия}
\begin{itemize}
  \item Участие в образовательных, социальных и благотворительных проектах.
  \item Организация мероприятий и работа с участниками.
\end{itemize}

% ===== КУРСЫ =====
\section{Курсы}

\cvevent{Яндекс Кружок}{Алгоритмы: параллели А и А'}{2022 -- 2024}{Москва, Россия}

% ===== ДОСТИЖЕНИЯ =====
\section{Достижения}

\begin{itemize}
  \item Призёр олимпиад по информатике и математике (2018--2025)
  \item Codeforces: рейтинг \textbf{1800+}
\end{itemize}

\end{document}
